\clearpage
\fancyhead[L]{\small JUNO.}
\hypertarget{native-265-lower}{}
\begin{center}
JUNO.

\rule{0.6\linewidth}{0.4pt}

Monatliche Correspondenz. Vol. XI. pp. 184–193. February 1805.

\rule{0.6\linewidth}{0.4pt}
\end{center}

—\quad—\quad—

As always, the most interesting and important contributions came from Dr. \textsc{Gauss}, who continues his work with steady diligence. So we begin our account of the latest observations and calculations concerning Juno with his contributions once again. In November 1804 this astronomer had waited in vain for good meridian observations of Juno that would enable him to improve its elements further. But he had received none by December, and the positions calculated from the third set of elements differed from the observed ones by a few minutes. He therefore based the following fourth set of elements on several observations of Juno, some made by himself with the ring micrometer and others received from Dr. \textsc{Olbers}.

\begin{center}
\emph{Fourth Set of Elements of Juno’s Orbit.}
\end{center}

\begin{center}
\begingroup\fontsize{9}{11}\selectfont\setlength{\tabcolsep}{3pt}\renewcommand{\arraystretch}{1.1}
\begin{tabular}{lr}
Epoch 1805 . . . . . . . . . & $42^\circ\ 41'\ 34''$ \\
Aphelion . . . . . . . . . & $233\ 23\ 47$ \\
Ascending node . . . . . . & $171\ 4\ 12$ \\
Inclination . . . . . . . . . . . & $13\ 4\ 9$ \\
Daily motion . . . . . . . . & $812''091$ \\
Logarithm of the semiaxis . . . . & $0.426935$ \\
Eccentricity . . . . . . . . . . & $0.256841$ \\
\end{tabular}\endgroup
\end{center}

\indent From these elements Dr. \textsc{Gauss} calculated its geocentric path for December 1804, January and February 1805. We already gave the ephemeris for the middle month in the preceding issue, so here we give only those for December and February [These ephemerides are omitted from this volume.] —\quad—\quad—

Here are the three observations Dr. \textsc{Gauss} made in December and January.

\begin{center}
\begingroup\fontsize{9}{11}\selectfont\setlength{\tabcolsep}{3pt}\renewcommand{\arraystretch}{1.1}
\begin{tabular}{l|r|r|r}
& Mean time & R.A. $\JunoSymbol$ & \shortstack{Declinatio\\australis}\\ \hline
1804, 29 December & $6^{\mathrm u}\ 53^{\mathrm m}\ 43^{\mathrm s}$ & $7^\circ\ 23'\ 11''$ & $7^\circ\ 28'\ 2''$ \\
30 & $6\ 19\ 50$ & $7\ 44\ 5$ & $7\ 19\ 57$ \\
1805, 5 January & $5\ 6\ 57.5$ & $9\ 59\ 42$ & $6\ 30\ 56$ \\
\end{tabular}\endgroup
\end{center}
\par\Needspace{11\baselineskip}
\SourceLeaf{266}\par
\indent He wrote to us that the last observation seemed sound in itself, but the comparison with a star from the Histoire céleste made it very doubtful. On 5 January Dr. \textsc{Gauss} compared Juno eight times with the eighth-magnitude star whose position is given on p. 135 of the Histoire céleste as follows.

\begin{center}
\begingroup\fontsize{10}{12}\selectfont\setlength{\tabcolsep}{3pt}\renewcommand{\arraystretch}{1.1}
\begin{tabular}{c|r|r}
& R.A. & Zenith distance\\ \hline
8 & $0^{\mathrm u}\ 40^{\mathrm m}\ 14^{\mathrm s}$ & $55^\circ\ 19'\ 15''$ \\
\end{tabular}\endgroup
\end{center}


\noindent Yet even by eye he thought that this star and three others above it formed a rhomboidal figure different from the one implied by the Histoire céleste. As Dr. \textsc{Gauss} wanted a more accurate position for the star, we took the first clear evening to determine it and found that his suspicion was fully justified. Our result for the southern declination differed from $6^\circ\,28'$. We obtained $6^\circ\,23'\,9''$. On that basis the declination observed for Juno on 5 January would be $6^\circ\,25'\,47''$. The fourth set of elements gives the following values.

\begin{center}
\begingroup\fontsize{10}{12}\selectfont\setlength{\tabcolsep}{3pt}\renewcommand{\arraystretch}{1.1}
\begin{tabular}{l|r|r}
& R.A. $\JunoSymbol$ & Declination\\ \hline
1804, 30 December & $7^\circ\ 44'\ 14''$ & $7^\circ\ 20'\ 16''$ \\
1805, 5 January & $9\ 59\ 44$ & $6\ 25\ 58$ \\
\end{tabular}\endgroup
\end{center}


\noindent After the latter reduction, the observed positions agree very well with these values. In determining the position of that doubtful star, however, we must point out that it was still so light at culmination that these observations can hardly count as fully reliable. —\quad—\quad—

[With this fourth set of elements, \textsc{Gauss} calculated the zodiac of Juno following the rules printed above on \hyperlink{native-111}{p. 106}, Mon. Corr. vol. XI, pp. 225–228. March 1805.]

\hfill [\textsc{von Zach.}]

